\chapter{پیشگفتار}
\begin{quote}
	{\lotus
		ریاضیات، زیباترین و قدرتمندترین مخلوق روح بشر است. ریاضیات به اندازه‌ی آدمی قدمت دارد. \quad - باناخ%
		\LTRfootnote{ Stefan Banach}
	}
\end{quote}
این جزوه را در مقام یک معلم نمی‌نویسم بلکه در مقام یک دانش‌آموز علاقه‌مند به ریاضی می‌نویسم. کتاب‌های ریاضی در نظام جدید آموزشی ایران در سطحی ساده قرار دارند و برای علاقه‌مندان جدی به ریاضی کافی نیست. درست است که بر اساس همان مطالب نیز می‌شود مسائل گوناگون زیادی را بررسی کرد، اما من به عنوان یک دانش‌آموز که دوست دارد بیشتر بداند، نیاز به منابع اضافی دارم. با خودم گفتم چرا به سراغ کتاب‌های نظام قدیم نروم؟ شاید به دلیل قدیمی بودن و کیفیت چاپ پایین‌تر نادیده گرفته شوند، اما از نظر ریاضی غنی و سرشار از مطالب مفید هستند. همچنین برای تکمیل برخی مباحث از کتاب‌های ریاضی خارج از ایران استفاده کردم تا با شیوه‌های دیگر آموزش ریاضی نیز آشنا شوم و بهره ببرم.

می‌دانم از منظر قانون کپی‌رایت کار درستی انجام نمی‌دهم، بنابراین سعی می‌کنم این جزوه را به صورت عمومی منتشر نکنم و فقط در صورت نیاز برای افراد نزدیک به اشتراک بگذارم، مثل یک جزوه‌ی درسی شخصی. صادقانه، مطالب این جزوه را از منابع مختلف دزدیده‌ام و گردآوری کرده‌ام! فهرست تمام منابع را در انتهای جزوه آورده‌ام و اگر تصویری نیز در جزوه استفاده کرده‌ام، منبعش را در قسمت منابع تصاویر ذکر کرده‌ام.

برای علاقه‌مندان به زبان انگلیسی، واژه‌نامه‌ای در انتهای جزوه قرار داده‌ام و معادل انگلیسی واژه‌های کلیدی هر بخش را نوشته‌ام. تلاش کنید زبان انگلیسی را یاد بگیرید و از آن فرار نکنید.

درباره‌ی نقل‌قول‌هایی که در آغاز هر فصل آورده شده است، اطمینان کامل ندارم. بیشتر آن‌ها از کتاب دکتر میرزاوزیری و همسرشان برداشته شده است و برخی نیز از منابع دیگر. امیدوارم معتبر و موثق باشند. اگر هم نباشند، همچنان نقل‌قول‌های خیالی جالبی هستند!

نقل‌قول‌های جالب و آموزنده‌ای از هالموس%
\LTRfootnote{ Paul Richard Halmos}
وجود دارد که دوست دارم در ادامه بیاورم: \\
- نباید صرفاً این را بخوانید، با آن مبارزه کنید! سؤالات خودتان را بپرسید، به مثال‌های خودتان بنگرید، برهان‌های خودتان را کشف نمایید. آیا فرض‌ها واقعاً لازمند؟ آیا عکس حکم نیز برقرار است؟ در حالت خاص متعارف چه اتفاقی می‌افتد؟ در مورد حالات استثنایی چطور؟ برهان، کجا از فرض استفاده می‌کند؟ \\
- انبوه مناسبی از مثال‌ها که تا حد امکان بزرگ باشند، برای درک عمیق هر مفهومی ضروری است و هرگاه که من می‌خواهم چیز جدیدی را یاد بگیرم، اولین کار من این است که نمونه‌ای از آن بسازم. \\
- تنها راه یاد گرفتن ریاضیات انجام دادن آن است.
\bigskip
\\
این جزوه در مرحله‌ی آماده‌سازی است و مشخص نیست در چه زمانی کامل می‌شود.
